\begin{titlepage}
\centering
{\Large Notas didácticas de una entrevista a fondo\par}
\vspace{1.0cm}
{\huge\bfseries YouWare, OS Agent y el nuevo paradigma de producto en el emprendimiento de aplicaciones de IA}\par
\vspace{0.5cm}
{\Large Zhang Xiaojun conversa con Ming Chaoping (明超平), fundador de YouWare\par}
\vspace{0.35cm}
{\large Elaborado a partir del video público, la transcripción, la estructura de capítulos, la descripción del video y diagramas conceptuales propios\par}
\vspace{0.3cm}
{\large 2025-05-29\par}
\vspace{0.85cm}
\includegraphics[width=0.88\textwidth,height=0.42\textheight,keepaspectratio]{cover.jpg}\par
\vfill
\begin{tcolorbox}[width=0.92\textwidth, colback=black!2!white, colframe=black!55, sharp corners]
\textbf{Título del video}: 101. Entrevista de 3 horas con Ming Chaoping, fundador de YouWare: hoy el Agent es como un gorila que acaba de tomar un atizador\par
\textbf{Canal del video}: Zhang Xiaojun Podcast\par
\textbf{Duración del video}: 02:42:45\par
\textbf{Enlace del video}: \href{\videourl}{\nolinkurl{\videourl}}\par
\textbf{Nota sobre los materiales}: YouTube no ofrece subtítulos descargables; el texto se elaboró a partir de una transcripción por bloques con faster-whisper large-v3 (CPU, int8), la descripción del video, la estructura de capítulos y figuras didácticas propias. Las tomas fijas de la entrevista no se reutilizan en el cuerpo del texto.\par
\textbf{Página del proyecto}: \href{\repourl}{\nolinkurl{\repourl}}\par
\end{tcolorbox}
\end{titlepage}

\begingroup
\small
\setlength{\parskip}{0pt}
\renewcommand{\baselinestretch}{0.92}\selectfont
\tableofcontents
\endgroup
\newpage

